\documentclass[12pt,a4paper]{article}
% Drake_preamble.tex - LaTeX Preamble Template
% 作者: 謝慕揚, MD, PhD, FESC
% 
% 使用說明:
% 1. 表格請使用 longtable，不要有垂直分隔線 |
% 2. 表格內換行使用 \newline，不要使用 \\
% 3. 藥物名稱、檢驗、檢查請維持英文
% 4. Reference 格式請使用 \href 加上驗證過的 DOI

% Including essential packages for Chinese typesetting
\usepackage{xeCJK}
\usepackage{fontspec}
% Configuring Chinese font with Noto Serif CJK TC, fallback to SimSun
\IfFontExistsTF{Noto Serif CJK TC}{
  \setCJKmainfont{Noto Serif CJK TC}
  \setCJKsansfont{Noto Serif CJK TC}
  \setCJKmonofont{Noto Serif CJK TC}
}{\setCJKmainfont{SimSun}
  \setCJKsansfont{SimSun}
  \setCJKmonofont{SimSun}
}

% Including other necessary packages
\usepackage{geometry}
\usepackage{titlesec}
\usepackage{enumitem}
\usepackage{xcolor}
\usepackage{colortbl}
\usepackage{tcolorbox}
\usepackage{booktabs}
\usepackage{longtable}
\usepackage{array}
\usepackage{graphicx}
\usepackage{tikz}
\usepackage{fancyhdr}
\usepackage{multicol}
\usepackage{datetime2}
\usepackage{hyperref}
\usepackage{mdframed}
\usepackage{amsmath}
\usepackage{amssymb}
\usepackage{multirow}

\hypersetup{
    colorlinks=true,
    linkcolor=emergency_blue,
    citecolor=emergency_blue,
    urlcolor=emergency_blue,
    pdfborder={0 0 0}
}

% Defining page geometry
\geometry{
  top=2.5cm,
  bottom=2.5cm,
  left=2cm,
  right=2cm,
  headheight=25pt
}

% Adjusting topmargin to compensate for increased headheight
\addtolength{\topmargin}{-10pt}

% Defining colors
\definecolor{emergency_red}{RGB}{186,24,27}
\definecolor{emergency_blue}{RGB}{0,114,188}
\definecolor{emergency_gray}{RGB}{120,120,120}
\definecolor{highlight_yellow}{RGB}{255,250,205}

% Setting title formats
\titleformat{\section}[block]{\Large\bfseries\color{emergency_red}}{\thesection}{10pt}{}
\titleformat{\subsection}[block]{\large\bfseries\color{emergency_blue}}{\thesubsection}{8pt}{}
\titleformat{\subsubsection}[block]{\normalsize\bfseries\color{emergency_gray}}{\thesubsubsection}{6pt}{}

% Defining custom environment for important notes
\newmdenv[
    linecolor=emergency_red,
    backgroundcolor=highlight_yellow,
    linewidth=2pt,
    topline=true,
    bottomline=true,
    leftline=true,
    rightline=true,
    innertopmargin=10pt,
    innerbottommargin=10pt,
    innerleftmargin=10pt,
    innerrightmargin=10pt
]{importantbox}

% Defining note command with tcolorbox
\newcommand{\note}[1]{
    \par
    \noindent
    \begin{tcolorbox}[
        colback=emergency_blue!5!white,
        colframe=emergency_blue,
        title=\textbf{重點提醒},
        fonttitle=\bfseries,
        boxrule=0.5pt,
        arc=3pt,
        left=6pt,
        right=6pt,
        top=6pt,
        bottom=6pt
    ]
    #1
    \end{tcolorbox}
}

% Key findings box
\newcommand{\keyfinding}[1]{
    \par
    \noindent
    \begin{tcolorbox}[
        colback=yellow!10!white,
        colframe=orange!75!black,
        title=\textbf{核心發現摘要},
        fonttitle=\bfseries,
        boxrule=0.5pt,
        arc=3pt
    ]
    #1
    \end{tcolorbox}
}

% Defining custom date format for ISO 8601
\DTMnewdatestyle{iso}{
  \renewcommand{\DTMdisplaydate}[4]{\number##1-\number##2-\number##3}
  \renewcommand{\DTMDisplaydate}[4]{\number##1-\number##2-\number##3}
}
\DTMsetdatestyle{iso}
\newcommand{\mydate}{\DTMtoday}


% Custom header for this document
\pagestyle{fancy}
\fancyhf{}
\fancyhead[L]{\textcolor{emergency_red}{Intensive Care Medicine 教學講義}}
\fancyhead[R]{\textcolor{emergency_blue}{VA-ECMO Echo Checklist}}
\fancyfoot[C]{\textcolor{emergency_gray}{\thepage}}
\renewcommand{\headrulewidth}{1pt}
\renewcommand{\headrule}{\hbox to\headwidth{\color{emergency_red}\leaders\hrule height \headrulewidth\hfill}}

\begin{document}

\begin{center}
{\LARGE\bfseries\textcolor{emergency_red}{Intensive Care Medicine Editorial}}\\[0.5cm]
{\Large\textbf{My Echo Checklist in Venoarterial ECMO Patients}}\\[0.3cm]
{\large VA-ECMO 病人的心臟超音波檢查清單}\\[0.5cm]
{\normalsize Intensive Care Med 2024;50:2158-2161. DOI: 10.1007/s00134-024-07659-2}\\[0.3cm]
{\normalsize 整理：謝慕揚 MD, PhD, FESC \quad 日期：\mydate}
\end{center}

\keyfinding{
\textbf{核心訊息：}Transthoracic echocardiography 是 VA-ECMO 管理的必要工具，本文提出系統性的九點檢查清單。

\textbf{實用重點：}
\begin{itemize}
    \item \textbf{每日評估：}LVEF、LVOT VTI、aortic valve opening、pericardial effusion
    \item \textbf{Weaning 標準：}LVOT VTI $\geq$12 cm、LVEF $\geq$20\%、S mitral wave $\geq$6 cm/s
    \item \textbf{LV unloading：}考慮 IABP 或 Impella 以預防 LV/aortic thrombosis
\end{itemize}
}

\tableofcontents
\newpage

%==============================================================
\section{文獻概述}
%==============================================================

本篇為 Intensive Care Medicine 的 Editorial，由法國 Sorbonne 大學 Pitié-Salpêtrière 醫院的 Saura 等人撰寫，提出 VA-ECMO 病人心臟超音波評估的系統性九點檢查清單。

\subsection{作者資訊}
\textbf{通訊作者：}Ouriel Saura, MD\\
\textbf{機構：}Service de Médecine Intensive-Réanimation, Institut de Cardiologie, Hôpital Pitié-Salpêtrière, Sorbonne Université, Paris, France

%==============================================================
\section{九點檢查清單總覽}
%==============================================================

\begin{longtable}{p{0.5cm}p{4cm}p{9.5cm}}
\toprule
\textbf{\#} & \textbf{檢查項目} & \textbf{重點內容} \\
\midrule
\endfirsthead

\toprule
\textbf{\#} & \textbf{檢查項目} & \textbf{重點內容} \\
\midrule
\endhead

\bottomrule
\endfoot

1 & Drainage cannula position & 最佳位置在 RA；若在 IVC 可能造成 suction \\
\midrule
2 & Drainage insufficiency & Subcostal view 評估；排除 tamponade、hypovolemia \\
\midrule
3 & ECMO flow 設定 & 根據 LVOT VTI、hemodynamic stability、lactate 調整 \\
\midrule
4 & Aortic valve opening & 預防 LV/aortic thrombosis；評估是否需要 unloading \\
\midrule
5 & Venting device 位置 & IABP 或 Impella 定位確認 \\
\midrule
6 & Pericardial effusion & 區分 true tamponade vs ECMO-induced pseudo-tamponade \\
\midrule
7 & Right-sided ECMO 設定 & 合併 LV support 時的 flow 調整 \\
\midrule
8 & Weaning 評估 & LVOT VTI $\geq$12 cm、LVEF $\geq$20\%、S wave $\geq$6 cm/s \\
\midrule
9 & Post-ECMO 併發症 & IVC/RA thrombosis、arterial pseudoaneurysm \\
\bottomrule
\end{longtable}

%==============================================================
\section{Point 1：Drainage Cannula 位置}
%==============================================================

\subsection{最佳位置}

\begin{itemize}
    \item \textbf{理想位置：}Drainage cannula tip 位於 \textbf{right atrium (RA)} 內
    \item 若位於 \textbf{IVC}：IVC 壁可能被 suction，阻礙血液引流
    \item 若位於 \textbf{SVC}：通常 ECMO flow 不受影響
\end{itemize}

\subsection{需注意的併發症}

\begin{itemize}
    \item Cannula tip 接觸 \textbf{interatrial septum}：
    \begin{itemize}
        \item 引流不足
        \item Endothelial injury
        \item Thrombus 形成風險
    \end{itemize}
    \item \textbf{Patent foramen ovale} 穿孔或重新開放：
    \begin{itemize}
        \item 含氧血從 LA 流向 return cannula
        \item 減少 LV preload 及 ejection
    \end{itemize}
\end{itemize}

\note{
\textbf{RA-PA double-lumen cannula：}
\begin{itemize}
    \item 血液 reinfuse 到 pulmonary artery
    \item 若 tip 位於 pulmonary valve 下方 → RV 支持不足
    \item 可能導致 RV dilation 及 tricuspid regurgitation
\end{itemize}
}

%==============================================================
\section{Point 2：Drainage Insufficiency 評估}
%==============================================================

\subsection{臨床表現}

\begin{itemize}
    \item ECMO lines \textbf{chattering}（管路震顫）
    \item ECMO flow 快速且短暫下降
    \item 原因：Multistage drainage cannula 吸附 IVC 或 RA 壁
\end{itemize}

\subsection{Subcostal View 評估重點}

\begin{enumerate}
    \item 確認 drainage cannula tip 位置（若在 IVC → 推進至 RA）
    \item 排除 significant pericardial effusion causing tamponade
    \item 確認 IVC 壁與 cannula 間有足夠空間（virtual space → 需要補液）
\end{enumerate}

\subsection{其他需排除的原因}

\begin{itemize}
    \item Pneumothorax
    \item Abdominal hypertension
\end{itemize}

\subsection{處理策略}

\begin{enumerate}
    \item 降低 pump RPM，然後緩慢增加
    \item 若仍不足 → 補液後再增加 RPM 及 ECMO flow
\end{enumerate}

%==============================================================
\section{Point 3：ECMO Flow 設定}
%==============================================================

\note{
\textbf{ECMO flow 設定考量因素：}
\begin{itemize}
    \item Residual stroke volume（LVOT VTI）
    \item Aortic valve opening 及 intracardiac thrombotic risk
    \item Mean arterial pressure
    \item Vasopressor 需求
    \item Lactate level/clearance
    \item 器官灌流臨床徵象（skin mottling、urine output、neurologic status）
\end{itemize}
}

\subsection{Flow 建議}

\begin{longtable}{p{4cm}p{10cm}}
\toprule
\textbf{Flow 設定} & \textbf{適用情況} \\
\midrule
\endfirsthead

\toprule
\textbf{Flow 設定} & \textbf{適用情況} \\
\midrule
\endhead

\bottomrule
\endfoot

\textbf{High flow (3-4 LPM)} &
• Severe hemodynamic instability\newline
• High vasoactive drugs doses\newline
• Elevated lactate\newline
• Poor residual ejection (LVOT VTI <5 cm) \\
\midrule
\textbf{Lower flow (2-3 LPM)} &
• Stabilized patients\newline
• 有 residual LV ejection \\
\bottomrule
\end{longtable}

\subsection{每日評估}

\begin{itemize}
    \item 從 ECMO 開始到 recovery，每日評估 \textbf{LVEF} 及 \textbf{LVOT VTI}
\end{itemize}

%==============================================================
\section{Point 4：Aortic Valve Opening}
%==============================================================

\subsection{問題背景}

\begin{itemize}
    \item Severe cases with very low residual cardiac ejection
    \item 特別是 high ECMO flow（decreased preload + increased afterload）
    \item Aortic valve 可能持續關閉
\end{itemize}

\subsection{併發症風險}

\begin{itemize}
    \item Blood stagnation
    \item \textbf{LV thrombus} 形成
    \item \textbf{Aortic root thrombus} 形成
\end{itemize}

\subsection{維持 Aortic Valve Opening 的策略}

\begin{enumerate}
    \item \textbf{減少 ECMO flow}（前提是 circulatory support 足夠）
    \item \textbf{Inotropic support}
    \item \textbf{LV unloading}（IABP）
    \item \textbf{Fluid challenge}
\end{enumerate}

\note{
\textbf{重要提醒：}
\begin{itemize}
    \item 有 atrioseptostomy 或 trans-septal LA drainage cannula 的病人，LV ejection 可能因 LV preload 明顯減少而惡化
    \item 若 aortic valve 持續關閉 → 考慮 direct venting（如 Impella）以預防 LV thrombus
\end{itemize}
}

%==============================================================
\section{Point 5：Venting Device 位置與設定}
%==============================================================

\subsection{IABP (Intra-Aortic Balloon Pump)}

\begin{longtable}{p{4cm}p{10cm}}
\toprule
\textbf{項目} & \textbf{內容} \\
\midrule
\endfirsthead

\toprule
\textbf{項目} & \textbf{內容} \\
\midrule
\endhead

\bottomrule
\endfoot

\textbf{作用機轉} &
• 降低 LV afterload\newline
• 促進 residual LV ejection\newline
• 促進 aortic valve opening\newline
• 降低 pulmonary edema 風險\newline
• 可能降低死亡率 \\
\midrule
\textbf{位置確認} & Suprasternal TTE view：tip 位於 left subclavian artery 起始處 \\
\bottomrule
\end{longtable}

\subsection{Impella}

\begin{longtable}{p{4cm}p{10cm}}
\toprule
\textbf{項目} & \textbf{內容} \\
\midrule
\endfirsthead

\toprule
\textbf{項目} & \textbf{內容} \\
\midrule
\endhead

\bottomrule
\endfoot

\textbf{作用機轉} &
• 降低 LV end-diastolic volume 及 pressure\newline
• 降低 pulmonary edema 風險\newline
• 降低 intracardiac/intra-aortic thrombus 風險 \\
\midrule
\textbf{正確位置} & Inlet 位於 aortic valve annulus 下方 \textbf{3-4 cm}，在 LV 內 \\
\midrule
\textbf{確認方法} & Color Doppler aliasing 確認血液 return 至 ascending aorta \\
\midrule
\textbf{位置不當} & 若 positioned too deeply → 血液 reinjected 回 LV\newline
→ 無法提供支持且增加 LV distension 風險 \\
\bottomrule
\end{longtable}

%==============================================================
\section{Point 6：Pericardial Effusion 與 Tamponade}
%==============================================================

\subsection{評估挑戰}

\begin{itemize}
    \item VA-ECMO 病人的 pericardial effusion 評估具挑戰性
    \item \textbf{RA collapse} 可能原因：
    \begin{itemize}
        \item Drainage cannula 的 suction
        \item Compressive pericardial effusion（真正 tamponade）
    \end{itemize}
\end{itemize}

\subsection{區分 True Tamponade vs Pseudo-Tamponade}

\begin{longtable}{p{4cm}p{10cm}}
\toprule
\textbf{徵象} & \textbf{解讀} \\
\midrule
\endfirsthead

\toprule
\textbf{徵象} & \textbf{解讀} \\
\midrule
\endhead

\bottomrule
\endfoot

IVC dilation + RV collapse & 支持 \textbf{true tamponade} → 需要 drainage \\
\midrule
Transient ECMO flow 減少後 tamponade appearance 消失 & 支持 \textbf{ECMO-induced pseudo-tamponade} \\
\bottomrule
\end{longtable}

%==============================================================
\section{Point 7：Right-Sided ECMO 合併 LV Support}
%==============================================================

\subsection{適用情況}

\begin{itemize}
    \item RA-PA double-lumen cannula 用於 LVAD 植入後嚴重 RV failure
    \item Peripheral VA-ECMO 相關併發症（cardiac akinesia、高 thrombosis 風險、嚴重 cannula site infection）轉換為 double central ECMO（RA-PA + LV-AA circuits）
\end{itemize}

\subsection{Right-Sided ECMO Flow 不足的後果}

\begin{enumerate}
    \item RV dilation
    \item Interventricular septum 向左移位
    \item LV support suction 導致 LV collapse
    \item Intraventricular obstruction
    \item Left-sided low-flow
\end{enumerate}

\subsection{處理策略}

\begin{itemize}
    \item 增加 right-sided ECMO flow（可能需要 fluid challenge）
    \item 目標：達到 interventricular septum neutral position
    \item 恢復 left-sided adequate flow
\end{itemize}

%==============================================================
\section{Point 8：Weaning 評估}
%==============================================================

\note{
\textbf{Weaning Trial 方法：}
\begin{itemize}
    \item 將 ECMO flow 短暫降至 \textbf{1 L/min}，持續 \textbf{5-10 分鐘}
    \item 結合 clinical hemodynamic assessment 與 Doppler echocardiographic evaluation
\end{itemize}
}

\subsection{Successful Weaning 預測指標}

\begin{longtable}{p{5cm}p{9cm}}
\toprule
\textbf{參數} & \textbf{Weaning 成功閾值} \\
\midrule
\endfirsthead

\toprule
\textbf{參數} & \textbf{Weaning 成功閾值} \\
\midrule
\endhead

\bottomrule
\endfoot

\textbf{LVOT VTI} & $\geq$12 cm \\
\midrule
\textbf{LVEF} & $\geq$20\% \\
\midrule
\textbf{S mitral wave} & $\geq$6 cm/s \\
\midrule
\textbf{Vasopressor/Inotrope} & 低劑量 \\
\midrule
\textbf{Hemodynamic stability} & 維持穩定 \\
\bottomrule
\end{longtable}

\subsection{其他預測方法}

\begin{itemize}
    \item Ventricular interdependence 評估
    \item RV systolic function 及 coupling to pulmonary circulation
\end{itemize}

%==============================================================
\section{Point 9：Post-ECMO 併發症}
%==============================================================

\subsection{血管併發症}

Peripheral ECMO insertion site 常見併發症：

\begin{longtable}{p{4cm}p{10cm}}
\toprule
\textbf{併發症} & \textbf{說明} \\
\midrule
\endfirsthead

\toprule
\textbf{併發症} & \textbf{說明} \\
\midrule
\endhead

\bottomrule
\endfoot

\textbf{IVC Thrombus} &
• 「Phantom of the cannula」外觀\newline
• 或附著於 IVC 壁的 thrombus\newline
• 有 pulmonary embolism 風險\newline
• 需要 anticoagulation \\
\midrule
\textbf{RA Thrombus} & 較少見 \\
\midrule
\textbf{Arterial Pseudoaneurysm} & Cannulation site；可用 echo-Doppler 偵測 \\
\midrule
\textbf{Stenosis} & Arterial cannulation site \\
\midrule
\textbf{Dissection} & Arterial cannulation site \\
\bottomrule
\end{longtable}

%==============================================================
\section{Take-Home Message}
%==============================================================

\begin{enumerate}
    \item \textcolor{red}{\textbf{每日評估：}}Doppler echocardiography 結合 clinical evaluation，是 VA-ECMO 病人 real-time management 的必要工具
    \item \textcolor{red}{\textbf{系統性檢查：}}使用九點檢查清單，systematic daily evaluation 各項 cardiac function 及 potential complications
    \item \textcolor{red}{\textbf{Weaning 標準：}}LVOT VTI $\geq$12 cm、LVEF $\geq$20\%、S mitral wave $\geq$6 cm/s，加上低劑量 vasopressors 及 hemodynamic stability
    \item \textcolor{red}{\textbf{Thrombosis 預防：}}維持 aortic valve opening；必要時使用 IABP 或 Impella 進行 LV unloading
\end{enumerate}

%==============================================================
\section{縮寫對照表}
%==============================================================

\begin{longtable}{p{3cm}p{11cm}}
\toprule
\textbf{縮寫} & \textbf{全名} \\
\midrule
\endfirsthead

\toprule
\textbf{縮寫} & \textbf{全名} \\
\midrule
\endhead

\bottomrule
\endfoot

VA-ECMO & Venoarterial Extracorporeal Membrane Oxygenation \\
TTE & Transthoracic Echocardiography \\
TEE & Transesophageal Echocardiography \\
RA & Right Atrium (右心房) \\
LA & Left Atrium (左心房) \\
LV & Left Ventricle (左心室) \\
RV & Right Ventricle (右心室) \\
IVC & Inferior Vena Cava (下腔靜脈) \\
SVC & Superior Vena Cava (上腔靜脈) \\
LVOT VTI & Left Ventricular Outflow Tract Velocity Time Integral \\
LVEF & Left Ventricular Ejection Fraction (左心室射血分數) \\
IABP & Intra-Aortic Balloon Pump (主動脈內氣球幫浦) \\
LVAD & Left Ventricular Assist Device (左心室輔助裝置) \\
LPM & Liters Per Minute \\
RPM & Revolutions Per Minute \\
PA & Pulmonary Artery (肺動脈) \\
AA & Aortic Artery (主動脈) \\
\bottomrule
\end{longtable}

%==============================================================
\section{參考文獻}
%==============================================================

\begin{enumerate}
    \item Saura O, Combes A, Hekimian G. My echo checklist in venoarterial ECMO patients. \href{https://doi.org/10.1007/s00134-024-07659-2}{\textit{Intensive Care Med}. 2024;50:2158-2161.}

    \item Douflé G, Dragoi L, Morales Castro D, et al. Head-to-toe bedside ultrasound for adult patients on extracorporeal membrane oxygenation. \href{https://doi.org/10.1007/s00134-024-07333-7}{\textit{Intensive Care Med}. 2024;50:632-645.}

    \item Ezad SM, Ryan M, Donker DW, et al. Unloading the left ventricle in venoarterial ECMO. In whom, when, and how? \href{https://doi.org/10.1161/CIRCULATIONAHA.122.062371}{\textit{Circulation}. 2023;147:1237-1250.}

    \item Aissaoui N, Luyt C-E, Leprince P, et al. Predictors of successful extracorporeal membrane oxygenation (ECMO) weaning after assistance for refractory cardiogenic shock. \href{https://doi.org/10.1007/s00134-011-2358-2}{\textit{Intensive Care Med}. 2011;37:1738-1745.}

    \item Lorusso R, Shekar K, MacLaren G, et al. ELSO interim guidelines for venoarterial extracorporeal membrane oxygenation in adult cardiac patients. \href{https://doi.org/10.1097/MAT.0000000000001510}{\textit{ASAIO J}. 2021;67:827-844.}
\end{enumerate}

\end{document}
